\documentclass[11pt,a4paper]{article}
\usepackage{fontspec,amsmath,amssymb,amsthm,microtype}
\setmainfont{lmroman10-regular.otf}[BoldFont=lmroman10-bold.otf,ItalicFont=lmroman10-italic.otf,BoldItalicFont=lmroman10-bolditalic.otf]
\setmonofont{lmmono10-regular.otf}[Scale=MatchLowercase]
\usepackage[margin=25mm]{geometry}
\usepackage[unicode,hidelinks]{hyperref}
\setlength{\emergencystretch}{3em}
\setlength{\parskip}{0.4em}
\setlength{\parindent}{0pt}
\allowdisplaybreaks[2]
\raggedbottom
\title{\texorpdfstring{An explicit quadratic bias envelope and second-\allowbreak{}order accuracy step}{An explicit quadratic bias envelope and second-order accuracy step}}
\author{LangevinLean mathematical exposition}
\date{}
\begin{document}\maketitle
\noindent Source module: \texttt{Kinetic\allowbreak{}Langevin/\allowbreak{}Second\allowbreak{}Order\allowbreak{}Budget.\allowbreak{}lean}. The accompanying HTML pairs this mathematical exposition with the exact Lean statements and records the scope of the formalization.


The stronger source Gaussian tail allows a completely explicit second-order accuracy prescription. The exponentially small term is bounded by a quadratic envelope, making the entire persistent floor at most \(B_p h^2\) for \(h\le1\). A positive minimum step then satisfies stability, the Gaussian cutoff dimension restriction, and the epsilon budget simultaneously.


\section{\texorpdfstring{Bound the genuine cutoff tail by a reciprocal}{Bound the genuine cutoff tail by a reciprocal}}


For \(x>0\), \(e^x\ge1+x\ge x\), hence \(e^{-x}\le1/x\). Apply this at \(x=m^2/(160M_2^2h^2)\) and use the previously proved sharper Gaussian tail. The result is \(e^{-(\alpha-p)/8}\le160M_2^2h^2/m^2\). This is a global inequality under the source restrictions, with no asymptotic qualification.



{\raggedright\noindent Lean declarations: \texttt{exp\_\allowbreak{}neg\_\allowbreak{}le\_\allowbreak{}reciprocal}, \texttt{second\allowbreak{}Order\allowbreak{}Cutoff\_\allowbreak{}tail\_\allowbreak{}quadratic}.\par}

\section{\texorpdfstring{Collect the complete floor into a quadratic coefficient}{Collect the complete floor into a quadratic coefficient}}


Write \(C_p=M_2\sqrt{p^2+2p}+M\sqrt{Mp}\). The exact floor is \(\sqrt2h^2C_p/m+(8(M-m)/m)h e^{-(\alpha-p)/8}\). Since \(M-m\le M\), the tail envelope bounds the second term by \(1280MM_2^2h^3/m^3\). For \(0<h\le1\), \(h^3\le h^2\). Thus the full floor is at most \(B_ph^2\), where \(B_p=\sqrt2 C_p/m+1280MM_2^2/m^3\). Every coefficient is explicit, and the exact fourth moment remains valid in dimension one.



{\raggedright\noindent Lean declarations: \texttt{second\allowbreak{}Order\allowbreak{}Bias\allowbreak{}Coefficient}, \texttt{second\allowbreak{}Order\_\allowbreak{}bias\_\allowbreak{}quadratic}.\par}

\section{\texorpdfstring{Choose a positive step satisfying every restriction}{Choose a positive step satisfying every restriction}}


For positive dimension and accuracy choose the minimum of \(1\), \(m/(5\gamma M)\), \(m/(4\sqrt{5p}M_2)\), and \(\sqrt{0.94\varepsilon/B_p}\). All four terms are positive. The second and third enforce the two source step restrictions; the last gives \(B_ph^2\le0.94\varepsilon\). Lean proves each minimum bound, compares nonnegative squares for the final term, and multiplies through positive denominators. For fixed curvature constants, \(B_p\) grows linearly with \(p\), so the displayed step and the logarithmic count exhibit the source \(\sqrt{p/\varepsilon}\) iteration dependence.



{\raggedright\noindent Lean declarations: \texttt{second\allowbreak{}Order\allowbreak{}Accuracy\allowbreak{}Step}, \texttt{second\allowbreak{}Order\allowbreak{}Accuracy\allowbreak{}Step\_\allowbreak{}spec}.\par}

\end{document}
